% ===========================================================================
%  metadati.tex — dati di frontespizio e identificazione della tesi
%
%  QUESTO È L'UNICO FILE DA COMPILARE A MANO PRIMA DELLA STAMPA.
%  Tutti i valori sono usati dal frontespizio (parti/frontespizio.tex) e dai
%  metadati PDF impostati da hyperref in preambolo.tex.
%
% ===========================================================================

% --- Istituzione -----------------------------------------------------------
\newcommand{\Ateneo}{Università Cattolica del Sacro Cuore}
\newcommand{\Dipartimento}{Dipartimento di Scienze Matematiche, Fisiche e Naturali}
\newcommand{\CorsoDiLaurea}{Corso di Laurea in Matematica a Curriculum Informatica}
\newcommand{\TipoLaurea}{Tesi di Laurea Triennale}
\newcommand{\AnnoAccademico}{2025/2026}
\newcommand{\LogoAteneo}{figure/logo_ateneo}   % senza estensione; commentare se assente

% --- Titolo ----------------------------------------------------------------
\newcommand{\TitoloTesi}{%
  Componenti quantistici nei transformer}
\newcommand{\SottotitoloTesi}{%
  uno studio sistematico su simulatori quantistici}

% --- Persone ---------------------------------------------------------------
\newcommand{\Candidato}{Lorenzo Maiuri}
\newcommand{\MatricolaCandidato}{5303839}
\newcommand{\Relatore}{Prof. Daniele Toti}
\newcommand{\Coautore}{Leonardo Pollini}            % autore del capitolo 2

% --- Parole chiave (metadati PDF) ------------------------------------------
\newcommand{\ParoleChiave}{%
  quantum machine learning, circuiti variazionali, Transformer,
  Vision Transformer, regolarizzazione implicita, risultati negativi}
